\section{Exact collisions with two independent target directions}
\label{sec:rank-two}
The preceding examples use geodesic images in the hyperbolic
target. A periodic Gowdy-to-Ernst transformation carries the
full-metric collision to a pair of solutions whose target maps
have differential of rank two on open subsets.
The transformation is standard \cite[\S 2, Eq.~(7)]{rw}.
Its action on this observation record follows from the common
current and metric sources of the pair.

\begin{theorem}
\label{th:rank-two}
For every $R>1$ there is a placement $b/a=R$ at which two
distinct vacuum Gowdy solutions with rank-two target
differential on open sets give the same induced spatial
metrics and the same $Q_t$ at both times, with different
$P_t(a)$. The measured $Q$ is spatially nonconstant at both
times. There is also a two-mode example for which no
time-orientation-preserving isometry that preserves the two
boundary slices and intertwines the torus actions can relate
the solutions, even if the Killing basis is allowed to change.
\end{theorem}

\begin{proof}
Use the geodesic-image pair of Theorem~\ref{th:exact-inhomogeneous}.
Its conserved target current is
$e^{2P}\,dQ/dw=q e^{P_0}=C$ on both branches.
Denote the transformed fields by tildes. In areal time,
\begin{equation}
 \widetilde P=-P-\log t,\quad
 \widetilde Q_t=Ct w_\theta,\quad
 \widetilde Q_\theta=Ct w_t,\quad
 \widetilde\lambda=\lambda-2P-\log t+C_0 .
 \label{eq:ernst}
\end{equation}
The last formula follows by substituting the first three into
both metric constraints. The same constant $C_0$ is used
on both branches.
The $\widetilde Q$ derivatives form an exact one-form on
the cylinder. Closedness follows from the scalar equation.
The angular period vanishes because $w$ has zero spatial mean.
Thus $\widetilde Q$ can be chosen equal on the two histories
for all times, with a common additive constant.
At both nodes $w=w_\theta=0$, so $\widetilde P$ and
$\widetilde\lambda$ agree and
$\widetilde Q_t=0$. At each node $w_t$ is nonzero as a
function of $\theta$, by uniqueness of the scalar ODE.
Consequently $\widetilde Q$ is nonconstant there.
At the initial slice
\[
 \det \frac{\partial(\widetilde P,\widetilde Q)}
                  {\partial(t,\theta)}
     =-C w_t(1+aP_w w_t),\qquad
 \widetilde P_t=-P_w w_t-1/a,\qquad P_w=\pm g.
\]
For a nonzero sufficiently small amplitude the determinant
does not vanish where $w_t\ne0$. This persists on open sets,
and the measured time velocities of $\widetilde P$ differ
between the branches.

To exclude boundary reparametrizations in one example,
write $\nu_j(R)^2$ for the $j$th radial Dirichlet eigenvalue
on $[1,R]$ used in \eqref{eq:invisible}.
The ratio $\nu_2(R)/\nu_1(R)$ varies continuously in $R$.
Rescaling the interval to length one and applying the
min--max principle gives a ratio tending to $2$ as $R\downarrow1$.
Rescaling the outer radius to one gives
$R\nu_j(R)\to j_{0,j}$ as $R\to\infty$, where $j_{0,j}$
is the $j$th positive zero of $J_0$.
To see the latter limit, extend annulus eigenfunctions by zero
in the hole for the lower eigenvalue bound. For the upper bound,
multiply the disk radial eigenfunctions by a cutoff that is
zero below $R^{-1}$ and rises linearly in log radius to one
at $R^{-1/2}$. Its squared radial gradient norm is
$2/\log R$, so the finite-dimensional disk trial space is
approximated in the weighted energy norm. The disk ratio
$j_{0,2}/j_{0,1}\simeq 2.295$ exceeds $13/6$.
Hence we can choose $a=\nu_1(R)/6$ and $b=Ra$ with
$13a=\nu_2(R)$. The modes of frequencies $6$ and $13$
then have simultaneous nodes. Take a nonzero sum of these
modes in $w$, choosing phases so that the common initial
$\widetilde Q$ after subtracting its mean is
\[
 F(\theta)=A_6\cos(6\theta)+A_{13}\cos(13\theta+\phi),
 \qquad A_6A_{13}\ne0,\quad \phi/\pi\notin\mathbb Q .
\]
All previous observations and the rank argument apply to
this phase, since they require only zero mean and both nodes.

On the initial slice the quotient circle metric is constant,
so an isometry compatible with the torus actions induces
$\theta\mapsto\pm\theta+c$ on that circle. A change of Killing
basis induces a hyperbolic isometry on the orbit shapes. Here
the initial orbit shapes range over an arc of a horocycle
of constant $\widetilde P$. A hyperbolic isometry that
carries this arc into the same horocycle therefore fixes
its ideal point and restricts to
$\widetilde Q\mapsto \pm\widetilde Q+\text{constant}$.
Taking means reduces the isometry condition on orbit shapes
to $F(\pm\theta+c)=\pm F(\theta)$.
Comparing the two Fourier coefficients shows that only the
identity on the quotient circle satisfies this condition.
The frequencies are coprime with opposite parity, so a
translation preserving both coefficients is the identity
and none flips both signs. A reflection would force
$\phi/\pi$ rational.
The two branches have different $\widetilde\lambda_t$
on the initial slice by \eqref{eq:ernst}.
The lapse agrees there and the determinant of the induced
metric is $t^{3/2}e^{-\widetilde\lambda/2}$.
Their mean curvatures thus differ as functions of this same
quotient circle. Preservation of the oriented initial
hypersurface would preserve mean curvature, a contradiction.
\end{proof}

The single-mode transformed examples approach a homogeneous
background of asymptotic speed one as the amplitude tends to
zero. Their parent solutions have uniformly subcritical
speed. In the transformed solutions the limiting speed
is $1-\sigma|\lim_{t\downarrow 0}t w_t|$ for sufficiently
small amplitude, reaching one where this limit vanishes.
This follows from the two ends of a nonvertical target
geodesic and from $t e^{\widetilde P}\widetilde Q_t
=Ct e^{-P}w_\theta\to0$. The rank-two result therefore concerns
finite-window recovery on smooth regular slices.
An illustrative two-mode computation gave $a=0.00173979$,
$b=0.453934$, and a maximum interior difference in
$\widetilde P$ of $0.10280$.

\begin{proposition}
\label{prop:ernst-fiber}
In the Gowdy-to-Ernst family just constructed, fix the common
initial metric, with a nonconstant $\widetilde Q(a)$, and
$\widetilde Q_t(a)=0$. If the underlying scalar phase has a
node at $b$, the original four observations leave a
one-parameter family of compatible histories. Adding the
terminal induced metric determines the magnitude of that
parameter and gives exactly two histories when the magnitude
is positive. One scalar reading of $\widetilde P$ at any
point with nonzero scalar phase then distinguishes the two.
\end{proposition}

\begin{proof}
Normalize the seed current by using the scalar
$\psi=Cw$ as the geodesic parameter.
The initial position of the seed is fixed in $P$ by
$\widetilde P(a)=-P_0-\log a$.
The seed $Q$ shift is immaterial to the transformed fields.
The initial $\widetilde Q$ determines $\psi_t(a)$ from
$a\psi_t(a)=\widetilde Q_\theta(a)$, with
$\psi(a)=0$. Thus $\psi$ is fixed on the whole interval.
The free parameter is $\zeta=P_\psi(0)$.
Current one gives $Q_\psi(0)=e^{-2P_0}$,
so the seed speed is
$\kappa=(\zeta^2+e^{-2P_0})^{1/2}$.
Its orbit value reads
\begin{equation}
 e^{-\widetilde P(t,\theta)-\log t-P_0}
  = \cosh(\kappa\psi)
       +\frac{\zeta}{\kappa}\sinh(\kappa\psi).
 \label{eq:ernst-fiber}
\end{equation}
At a node this is one for every $\zeta$, so the partial
record leaves the full family.

At a terminal node the metric difference is
\[
 \widetilde\lambda(b)-\widetilde\lambda(a)
   =-\kappa^2 I-\log(b/a),\qquad
 I=\int_a^b t(\psi_t^2+\psi_\theta^2)\,dt .
\]
The integral is constant on the quotient circle because
its angular derivative is $2[t\psi_t\psi_\theta]_a^b=0$.
It is positive by spatial averaging and the nontrivial
initial $\widetilde Q$. The terminal $\lambda$ therefore
determines $\kappa$, leaving precisely the two signs of
$\zeta$ when $\zeta\ne0$. Equation~\eqref{eq:ernst-fiber}
is affine in $\zeta$ at fixed $\kappa$, with a nonzero
coefficient wherever $\psi\ne0$. A scalar reading there
gives the sign and the value.
\end{proof}
